\section{Restriction to the open subbundle and faithfulness}\label{sec_TX_incl}

\subsection{The open embedding}\label{ss:inclusion_TX_Xdelta_overview}

Let $\CB=G/B$ denote the variety of complete flags
$F_\bullet=(0=F_0\subset F_1\subset\cdots\subset F_n=\C^n)$ in $\C^n$, and let
$\CN\subset\mathfrak{gl}_n(\C)$ be the nilpotent cone.
For a root ideal $\Psi=\Psi_{\rfunc}$, set
\[
  \Fn(\Psi)=\bigoplus_{\alpha\in\Psi}\mathfrak g_\alpha,
  \qquad T^*_\Psi\CB=G\times^B\Fn(\Psi).
\]
Via the standard trace pairing identifying $T^*\CB$ with $G\times^B\Fn(\Delta^+)$,
the vector bundle $T^*_\Psi\CB$ embeds into $T^*\CB$ as the subbundle
\[
  T^*_\Psi\CB
  =\{(F_\bullet,\xi)\in\CB\times\CN
        \mid \xi F_i\subset F_{\rfunc(i)}\text{ for all }1\leq i\leq n\}.
\]
In particular, $T^*_{\Delta^+}\CB=T^*\CB$.
The group $G\times S$ acts on $T^*_\Psi\CB$ by
\[
  (g,a)\cdot(F_\bullet,\xi)=(gF_\bullet,a^{-1}g\xi g^{-1}),
\]
where loop rotation scales the cotangent fibers with weight $-1$.

We use the classical lattice chart of
\cite{lusztig1981green,ngo1999faisceaux,mirkovic2003quiver}, as described for
$T^*\CB$ in \cite[Lemma~11.2]{cautis2018quantum}, and restrict it to the
subbundle $T^*_\Psi\CB$ inside Kato's compactification \cite{kato2023geometric}.

Recall from Section~\ref{sec_xpsi} that points of $\SX_\Psi$ are lattice flags $L_\bullet$ with
$L_0=t\CO^n$.
Set $V=(\CK/t\CO)^n$. Each flag $L_\bullet$ is uniquely determined by the flag of
finite-dimensional subspaces $M_i=L_i/L_0\subset V$.
Because $tM_i\subset M_{\rfunc(i)}\subset M_{i-1}$, each $M_i$ is contained in $t^{1-n}\CO^n/t\CO^n$.
With respect to the decomposition $V=(\CK/t\CO)^n=\bigoplus_{h\geq 0}t^{-h}\C^n$, let $p:V\to\C^n$ be the projection onto the degree-zero component $\CO^n/t\CO^n\cong\C^n$.
For any $\xi\in\CN$, define the linear map
\[
  q_\xi:\C^n\longrightarrow V,\qquad
  q_\xi(v)=\sum_{h=0}^{n-1}t^{-h}\xi^h v.
\]
Because $\xi^n=0$, multiplication by $t$ on $V$ satisfies $t q_\xi=q_\xi\xi$, while $p\circ q_\xi=\id_{\C^n}$.

\begin{proposition}\label{prop:T^*X_incl}
Under the identification of each lattice flag $L_\bullet\in\SX_\Psi$
with its quotient flag $M_\bullet=(L_i/L_0)_{i=0}^n$ in $V$, the map
\[
  \iota_\Psi:T^*_\Psi\CB\longrightarrow\SX_\Psi,\qquad
  (F_\bullet,\xi)\longmapsto\bigl(q_\xi(F_i)\bigr)_{i=0}^n
\]
is a $G\times S$-equivariant open immersion.
Its image consists of those lattice flags for which the restriction
$p|_{M_n}:M_n\to\C^n$ is an isomorphism.
\end{proposition}

\begin{proof}
The identities $tq_\xi=q_\xi\xi$ and $p\circ q_\xi=\id_{\C^n}$ give
$tq_\xi(F_i)\subset q_\xi(F_{\rfunc(i)})$.
On the stated open locus, the inverse morphism is given by
\[
 F_i=p(M_i),\qquad \xi=(p|_{M_n})\circ(t|_{M_n})\circ(p|_{M_n})^{-1},
\]
where $t|_{M_n}:M_n\to M_n$ denotes multiplication by $t$.
Indeed, $t^nM_n=0$ and every $v\in M_n$ satisfies
$v=\sum_{h=0}^{n-1}t^{-h}p(t^hv)=q_\xi(p(v))$.
In particular, $M_i=q_\xi(F_i)$ for all $i$, and $\xi F_i=p(tM_i)\subset p(M_{\rfunc(i)})=F_{\rfunc(i)}$, so $(F_\bullet,\xi)\in T^*_\Psi\CB$.
Both assignments are algebraic and $G$-equivariant.
The identity $a\cdot q_\xi(v)=q_{a^{-1}\xi}(v)$ for $a\in S$ proves $S$-equivariance.
\end{proof}


\subsection{Restriction of the operators}
\label{ss:inclusion_TX_Xdelta}

For $K=T\times S$ or $G\times S$, we also write $x_i$ for the restriction of
the tautological class to $T^*_\Psi\CB$; on the zero section it is
$-c_1^K(F_i/F_{i-1})$.
The bundle projection $T^*_\Psi\CB\to\CB$ provides a $G\times S$-equivariant deformation retraction
onto the zero section $\CB=G/B$.
Change of groups for the homogeneous space $G/B$ identifies its equivariant cohomology:
\[
 H_{G\times S}^*(T^*_\Psi\CB;\Z)
 \cong H_{B\times S}^*(\pt;\Z)
 \cong\Z[x_1,\dots,x_n,\epsilon],
\]
where the Chern classes $c_j\in\Lambda_{G\times S}$ map to $(-1)^j e_j(x_1,\dots,x_n)$.
For $K=T\times S$ or $G\times S$, let $J_K\subset R_K$ denote the ideal
\begin{align*}
 J_{T\times S}&=\bigl(e_j(x_1,\dots,x_n)-(-1)^je_j(t_1,\dots,t_n)\mid1\leq j\leq n\bigr),\\
 J_{G\times S}&=\bigl(e_j(x_1,\dots,x_n)-(-1)^jc_j\mid1\leq j\leq n\bigr).
\end{align*}
Together with Lemma~\ref{lem:borel}, this gives the canonical presentation
\[
 R_K/J_K\xrightarrow{\sim}H_K^*(T^*_\Psi\CB;\Z),\qquad
 x_i+J_K\longmapsto x_i.
\]
Thus the flag class $x_i$ is represented by the residue class $x_i+J_K$.

\begin{proposition}\label{prop:inclusion_TX_Xdelta_deg_hecke}
For either group $K$, restriction along $\iota_\Psi$ is the surjective
$\Hdeg'$-module homomorphism
\[
 R_K/I_{\rfunc,K}\longrightarrow R_K/J_K,\qquad
 x_i+I_{\rfunc,K}\longmapsto x_i+J_K.
\]
The restricted geometric correspondences are proper over both factors
and act by the operators $\BS_i$ on $R_K/J_K$.
\end{proposition}

\begin{proof}
On the zero section, $q_0(F_i)=F_i$, so the tautological lines restrict to the flag lines.
These classes generate the target over $\Lambda_K$, giving surjectivity
and $I_{\rfunc,K}\subset J_K$. Since $J_K$ is generated by polynomials symmetric in $x_1,\dots,x_n$,
it is preserved by each $\BS_i$, and Theorem~\ref{thm:geometric_hecke}
gives $\Hdeg'$-module compatibility.

Write $X=\SX_\Psi$ and identify $T^*_\Psi\CB$ with its image under $\iota_\Psi$.
For each $s_i\in S'$, every pair $(L_\bullet,L'_\bullet)\in C_i$ has the same terminal lattice $L_n=L'_n$,
and by Proposition~\ref{prop:T^*X_incl} the open condition that $p|_{L_n/L_0}$ is an isomorphism
depends only on $L_n$.
Hence
\[
 \operatorname{pr}_1^{-1}(T^*_\Psi\CB)\cap C_i
 =\operatorname{pr}_2^{-1}(T^*_\Psi\CB)\cap C_i
 =(T^*_\Psi\CB\times T^*_\Psi\CB)\cap C_i,
\]
so the projections $\operatorname{pr}_1|_{C_i},\operatorname{pr}_2|_{C_i}:C_i\to X$ form Cartesian squares over the open immersion $T^*_\Psi\CB\hookrightarrow X$, and their restrictions $C_i\cap(T^*_\Psi\CB\times T^*_\Psi\CB)\to T^*_\Psi\CB$ remain proper.
Proper base change therefore shows that restriction along $\iota_\Psi$ intertwines the action of $\Gamma_i$ with that of its restriction to $T^*_\Psi\CB\times T^*_\Psi\CB$ (when $\epsilon_i-\epsilon_{i+1}\notin\Psi$, the factor $\operatorname{pr}_1^*e_K(\mathcal E_i)$ restricts simultaneously with $[C_i]$); since $\iota_\Psi^*$ is surjective, the restricted correspondences act by $\BS_i$ on $R_K/J_K$.
\end{proof}

\begin{corollary}\label{cor:faithful_hecke}
For either group $K$, the $\Hdeg'$-action on $H_K^*(\SX_{\rfunc};\Z)$ is faithful over $\Z[\epsilon]$.
\end{corollary}

\begin{proof}
For $G\times S$, Proposition~\ref{prop:inclusion_TX_Xdelta_deg_hecke} realizes the faithful polynomial
representation of $\Hdeg'$ on $\Z[x_1,\dots,x_n,\epsilon]$ as a quotient module of $H_{G\times S}^*(\SX_{\rfunc};\Z)$, so Lemma~\ref{lem:polynomial_faithful} proves
the assertion. For $T\times S$, faithfulness follows from the equivariant injection
$H_{G\times S}^*(\SX_{\rfunc};\Z)\hookrightarrow H_{T\times S}^*(\SX_{\rfunc};\Z)$
of Lemma~\ref{lem:borel}.
\end{proof}

\begin{remark}
Faithfulness is taken over $\Z[\epsilon]$ rather than over $\Lambda_K$. For example, on the flag variety $\SX_\varnothing=\CB$, the element $e_1(x_1,\dots,x_n)+c_1$ acts by zero on
$H_{G\times S}^*(\CB;\Z)$, even though it is a nonzero element of $\Lambda_{G\times S}\otimes_{\Z[\epsilon]}\Hdeg'$.
\end{remark}

\begin{corollary}\label{thm:faithful_convolution_embedding}
For either group $K$ and $X=\SX_\Psi$, there is an injective homomorphism of graded $\Z[\epsilon]$-algebras
\[
 \Hdeg'\hookrightarrow\mathcal C_X^\bullet,\qquad
 x_j\longmapsto\Delta_{X*}x_j,\quad
 s_i\longmapsto\Gamma_i\quad(s_i\in S'),\quad
 \epsilon\longmapsto\epsilon[\Delta_X],
\]
where $\Delta_X:X\hookrightarrow X\times X$ is the diagonal embedding and $\Gamma_i$ are the correspondence classes of Section~\ref{sec:geometric_hecke}.
\end{corollary}

\begin{proof}
Under \eqref{eq:ambient_endomorphisms}, these classes act by the polynomial generators
and $\mathcal S_i$. Theorem~\ref{thm:geometric_hecke} and
Corollary~\ref{cor:faithful_hecke} therefore give the homomorphism and its injectivity.
\end{proof}

\subsection{Comparison with Steinberg correspondences}\label{sec:steinberg}

Assume now that $\Psi=\Delta^+$, and set $X=\SX_{\Delta^+}$, $M=T^*\CB$,
and $d=\dim_\C X=n(n-1)$.
After the lattice shift in Remark~\ref{rem:lattice_normalization},
$X$ is the convolution variety $Y(1^n)$ of \cite[Lemma~6.1]{cautis2018quantum}.
Write $Z=M\times_{\CN}M$ for the Steinberg variety and set
\[
 \overline Z=\{(L_\bullet,L'_\bullet)\in X\times X\mid L_n=L'_n\}
\]
with its reduced subscheme structure.
The open embedding of Proposition~\ref{prop:T^*X_incl} identifies
$Z$ with $\overline Z\cap(M\times M)$, as in \cite[Lemma~11.2]{cautis2018quantum}.

For $1\leq i<n$, let $\rho_i:\CB\to\CB_i$ forget the $i$-th subspace and put
$M_i=\rho_i^*T^*\CB_i\subset M$.
The classical simple Steinberg component is
\[
 Z_i=M_i\times_{T^*\CB_i}M_i
\]
(cf.\ \cite[\S12]{ginzburg1998geometric}).
Let $Y_i$ and $\pi_i:Y_i\to B_i$ be the divisor and forgetful morphism
constructed in Section~\ref{sec:geometric_hecke}.

\begin{proposition}\label{prop:compact_steinberg_hecke}
The closure of $Z_i$ in $X\times X$ is
\begin{equation}\label{eq:compact_simple_component}
 \overline Z_i=Y_i\times_{B_i}Y_i.
\end{equation}
It is smooth and irreducible of dimension $d$.
\end{proposition}

\begin{proof}
The variety $Y_i=\SX_{\rfunc_i^-}$ is an iterated projective bundle over a point,
and $Y_i\times_{B_i}Y_i$ is a $\PP^1$-bundle over $Y_i$.
It is therefore smooth and irreducible of dimension $d$.
It lies in $\overline Z$, since both flags have the same terminal lattice.

On $M$, the condition $tL_{i+1}\subset L_{i-1}$ becomes
$\xi F_{i+1}\subset F_{i-1}$, so $Y_i\cap M=M_i$.
By the inverse chart of Proposition~\ref{prop:T^*X_incl},
equality of the retained lattices is equivalent to equality of $\xi$ and the retained flag subspaces.
Thus the intersection of $Y_i\times_{B_i}Y_i$ with $M\times M$ is $Z_i$.
It is a nonempty open subset (containing pairs of equal flags with $\xi=0$),
hence dense in the irreducible fiber product, proving the closure assertion.
\end{proof}

Consequently, the generator images in Corollary~\ref{thm:faithful_convolution_embedding}
are $\Gamma_i=[\Delta_X]+[\overline Z_i]$.
On $M$, the class $\Gamma_i$ restricts to $[\Delta_M]+[Z_i]$,
whose action is identified with $\BS_i$ in Proposition~\ref{prop:inclusion_TX_Xdelta_deg_hecke}.
The corresponding $K$-theoretic construction and open comparison appear in
\cite[Lemma~9.3 and Proposition~11.3]{cautis2018quantum}.

